% =============================================================================
% Chapter 15: Scaling and Efficiency
% =============================================================================
\chapter{Scaling and Efficiency}
\label{ch:scaling-efficiency}

\begin{objectives}
    \item Understand scaling laws for language models
    \item Calculate compute requirements
    \item Optimize training efficiency
    \item Know when small models are sufficient
\end{objectives}

% -----------------------------------------------------------------------------
\section{Scaling Laws}
% -----------------------------------------------------------------------------

\subsection{The Chinchilla Insight}

Research by DeepMind (Hoffmann et al., 2022) established that:

\begin{equation}
\text{Optimal training} \implies N \propto D
\end{equation}

where $N$ is parameter count and $D$ is training data tokens.

\begin{keypoint}
For compute-optimal training, \textbf{scale data and parameters equally}. A 10M parameter model should train on $\sim$200M tokens; a 1B parameter model on $\sim$20B tokens.
\end{keypoint}

\subsection{Loss Scaling}

Language model loss follows a power law:
\begin{equation}
L(N, D) \approx \left(\frac{N_c}{N}\right)^{\alpha_N} + \left(\frac{D_c}{D}\right)^{\alpha_D} + L_\infty
\end{equation}

\begin{itemize}
    \item More parameters $\to$ lower loss
    \item More data $\to$ lower loss
    \item Diminishing returns in both directions
\end{itemize}

% -----------------------------------------------------------------------------
\section{Compute Requirements}
% -----------------------------------------------------------------------------

\subsection{FLOPs per Token}

For a transformer with $N$ parameters:
\begin{equation}
\text{FLOPs per token} \approx 6N \text{ (forward + backward)}
\end{equation}

\subsection{TryLM Compute}

For \TryLM ($N = 10.8$M, batch=256, steps=10K):
\begin{align*}
\text{Tokens} &= 256 \times 512 \times 10000 = 1.3\text{B} \\
\text{FLOPs} &= 6 \times 10.8\text{M} \times 1.3\text{B} = 84 \text{ TFLOPs}
\end{align*}

On an RTX 3080 ($\sim$30 TFLOPS fp16):
\[
\text{Time} \approx \frac{84 \text{ TFLOPs}}{30 \text{ TFLOPS}} \approx 45 \text{ minutes}
\]

% -----------------------------------------------------------------------------
\section{Efficiency Optimizations}
% -----------------------------------------------------------------------------

\subsection{torch.compile}

\begin{lstlisting}[style=python]
model = GPT(config)
model = torch.compile(model)  # 1.5-2x speedup
\end{lstlisting}

\texttt{torch.compile} fuses operations and optimizes GPU usage.

\subsection{Mixed Precision}

Float16 computation is 2x faster and uses half the memory.

\subsection{Flash Attention}

For longer sequences, Flash Attention provides $O(n)$ memory instead of $O(n^2)$:
\begin{lstlisting}[style=python]
from torch.nn.functional import scaled_dot_product_attention

# PyTorch 2.0+ automatically uses Flash Attention when available
attn = scaled_dot_product_attention(q, k, v, is_causal=True)
\end{lstlisting}

% -----------------------------------------------------------------------------
\section{When Small Models Suffice}
% -----------------------------------------------------------------------------

\TryLM shows that small models can be effective for specific domains:

\begin{table}[H]
\centering
\begin{tabular}{lll}
\toprule
\textbf{Task} & \textbf{Model Size} & \textbf{Notes} \\
\midrule
Simple text completion & 10--100M & Domain-specific \\
Code completion & 100M--1B & Syntax-focused \\
General knowledge & 7B+ & Broad training data \\
Complex reasoning & 70B+ & Emergent capabilities \\
\bottomrule
\end{tabular}
\caption{Model size requirements by task}
\end{table}

\begin{keypoint}
\textbf{Small models excel when}:
\begin{enumerate}
    \item The domain is narrow (children's stories)
    \item The task is well-defined (complete a sentence)
    \item Latency matters (real-time generation)
    \item Resources are limited (edge devices)
\end{enumerate}
\end{keypoint}

% -----------------------------------------------------------------------------
\section{Summary}
% -----------------------------------------------------------------------------

\begin{summary}
    \item \textbf{Scaling laws}: $L \propto N^{-\alpha} + D^{-\beta}$
    \item \textbf{Chinchilla optimal}: Train data tokens $\approx 20 \times$ parameters
    \item \textbf{FLOPs}: $\approx 6N$ per token (forward + backward)
    \item \textbf{Optimizations}: torch.compile, mixed precision, Flash Attention
    \item \textbf{Small models work} for narrow domains and specific tasks
\end{summary}
